% ----- CHAUPAI 40 -----

\newpage

\subsection*{\deva{चत्वारिंशत्तम चौपाई} (Fortieth Chaupai)}
\addcontentsline{toc}{subsection}{\deva{चत्वारिंशत्तम चौपाई} (Fortieth Chaupai) :  \deva{तुलसीदास सदा...}}

\begin{minipage}[t]{0.55\textwidth}
\vspace{0pt}

{\large \linespread{1.5}\selectfont
\deva{तुलसीदास सदा हरि चेरा।}\\
\deva{कीजै नाथ हृदय महँ डेरा॥}
\par}

\vspace{0.5em}

{\large \linespread{1.5}\selectfont
\telu{తులసీదాస సదా హరి చేరా।}\\
\telu{కీజై నాథ హృదయ మహఁ డేరా॥}
\par}

\vspace{0.5em}

{\large \linespread{1.3}\selectfont
\textit{tulasīdāsa sadā hari cerā |}\\
\textit{kījai nātha hṛdaya mahã ḍerā ||}
\par}
\end{minipage}%
\hfill
\begin{minipage}[t]{0.42\textwidth}
\vspace{0pt}
\begin{tcolorbox}[anchorbox]
\textbf{The Humble Heart:} Instead of demanding massive recognition or building monuments to ego, true greatness and peace begin with an unassuming, hospitable, and humble mind (\textit{Hriday Mahan Dera}).
\vspace{0.5em}\newline True greatness isn't found in massive monuments to the ego, but in building an unassuming, hospitable, and humble mind.\newline\null\hfill\href{https://www.valmiki.iitk.ac.in/sloka?field_kanda_tid=6&language=dv&field_sarga_value=128}{\textit{Yuddha Kanda, 128:83}}
\end{tcolorbox}

\end{minipage}

\vspace{0.5em}
\subsubsection*{\textbf{Visual Rhythm Map:}}
\begin{table}[h!]
\centering
\renewcommand{\arraystretch}{1.2}
\setlength{\tabcolsep}{0pt}
\begin{tabularx}{\textwidth}{|*{16}{>{\centering\arraybackslash\hsize=1.0\hsize}X|}}
\hline
\multicolumn{7}{|>{\centering\arraybackslash\hsize=7.0\hsize}X|}{\textbf{\laghu{तु}\laghu{ल}\guru{सी}\guru{दा}\laghu{स}} \par (7)} &
\multicolumn{3}{>{\centering\arraybackslash\hsize=3.0\hsize}X|}{\textbf{\laghu{स}\guru{दा}} \par (3)} &
\multicolumn{2}{>{\centering\arraybackslash\hsize=2.0\hsize}X|}{\textbf{\laghu{ह}\laghu{रि}} \par (2)} &
\multicolumn{4}{>{\centering\arraybackslash\hsize=4.0\hsize}X|}{\textbf{\guru{चे}\guru{रा}} \par (4)} \\
\hline
\end{tabularx}
\vspace{-1.5em}
\end{table}

\begin{table}[h!]
\centering
\renewcommand{\arraystretch}{1.2}
\setlength{\tabcolsep}{0pt}
\begin{tabularx}{\textwidth}{|*{16}{>{\centering\arraybackslash\hsize=1.0\hsize}X|}}
\hline
\multicolumn{4}{|>{\centering\arraybackslash\hsize=4.0\hsize}X|}{\textbf{\guru{की}\guru{जै}} \par (4)} &
\multicolumn{3}{>{\centering\arraybackslash\hsize=3.0\hsize}X|}{\textbf{\guru{ना}\laghu{थ}} \par (3)} &
\multicolumn{3}{>{\centering\arraybackslash\hsize=3.0\hsize}X|}{\textbf{\laghu{हृ}\laghu{द}\laghu{य}} \par (3)} &
\multicolumn{2}{>{\centering\arraybackslash\hsize=2.0\hsize}X|}{\textbf{\laghu{म}\laghu{हँ}} \par (2)} &
\multicolumn{4}{>{\centering\arraybackslash\hsize=4.0\hsize}X|}{\textbf{\guru{डे}\guru{रा}} \par (4)} \\
\hline
\end{tabularx}
\vspace{-1.5em}
\end{table}

\vspace{1em}
\textbf{ \deva{सरल-संस्कृतम्}:}\\
 \deva{हे नाथ! तुलसीदासः सर्वदा भगवतः हरिः (श्रीरामस्य) परिचारकः (सेवकः) अस्ति।}\\
 \deva{अतः हे स्वामिन्! भवान् मम हृदये एव स्वकीयम् आवासं (डेरा) करोतु॥}\\


Tulsidas is, and will always remain, the servant of Lord Hari (God).\\
Therefore, O Lord! Please make my heart your permanent dwelling.


\vspace{1em}
\newpage
\textbf{\deva{प्रतिपदार्थः} (Meaning of each word):}
\begin{table}[h!]
\centering
\renewcommand{\arraystretch}{1.2}
\begin{tabularx}{\textwidth}{@{} l l X X @{}}
\toprule
\textbf{Awadhi} & \textbf{Sanskrit} & \textbf{English} & \textbf{Grammar \& Notes} \\
\midrule
\deva{तुलसीदास} / \telu{తులసీదాస}  &  \deva{तुलसीदासः}  &  Tulsidas (The Author)  &   \\
\deva{सदा} / \telu{సదా}  &  \deva{सर्वदा}  &  Always  &   \\
\deva{हरि चेरा} / \telu{హరి చేరా}  &  \deva{हरेः सेवकः}  &  Servant of God  & `Chera` is pure Awadhi for servant \\
\deva{कीजै} / \telu{కీజై}  &  \deva{कुर्यात् / करोतु}  &  Please do / Make  & Respectful Imperative \\
\deva{नाथ} / \telu{నాథ}  &  \deva{नाथ}  &  O Lord  &   \\
\deva{हृदय महँ} / \telu{హృదయ మహఁ}  &  \deva{हृदय-मध्ये}  &  In the heart  & Locative postposition `mahan` \\
\deva{डेरा} / \telu{డేరా}  &  \deva{आवासम्}  &  Camp / Dwelling  &   \\
\bottomrule
\end{tabularx}
\end{table}
\vspace{0.5em}
\begin{tcolorbox}[poetrybox]
\textbf{Praasa (The Signature Verse):}\\
It is customary in Indian classical poetry for the author to gracefully 'sign' their work in the final verse. Here, the great classical poet Tulsidas identifies himself simply as a proud servant (\deva{चेरा}) of Hari. 
\end{tcolorbox}
\newpage
